\documentclass[12pt,a4paper]{article}

% Template de relatório — Cálculo Numérico
% Compilador recomendado no Overleaf: pdfLaTeX

\usepackage[utf8]{inputenc}
\usepackage[T1]{fontenc}
\usepackage[brazil]{babel}
\usepackage{lmodern}
\usepackage{geometry}
\geometry{top=2.5cm,bottom=2.5cm,left=3cm,right=2cm}
\usepackage{setspace}
\onehalfspacing
\usepackage{graphicx}
\usepackage{booktabs}
\usepackage{amsmath,amssymb}
\usepackage{float}
\usepackage{enumitem}
\usepackage{hyperref}
\usepackage{xcolor}
\usepackage{listings}
\usepackage[alf]{abntex2cite}

\hypersetup{
    colorlinks=true,
    linkcolor=black,
    citecolor=black,
    urlcolor=blue,
    pdftitle={Relatório — Cálculo Numérico},
    pdfauthor={UNIFESP}
}

\lstset{
    basicstyle=\ttfamily\small,
    breaklines=true,
    frame=single,
    numbers=left,
    numberstyle=\tiny,
    showstringspaces=false,
    captionpos=b
}

% Edite os campos abaixo
\newcommand{\TituloRelatorio}{Título do trabalho}
\newcommand{\Disciplina}{Cálculo Numérico}
\newcommand{\Universidade}{Universidade Federal de São Paulo (UNIFESP)}
\newcommand{\Curso}{Nome do curso}
\newcommand{\Ano}{2026}

\begin{document}

% -------------------- CAPA -----------------------
\begin{titlepage}
    \centering
    {\large \Universidade\par}
    \vspace{0.5cm}
    {\large \Curso\par}
    \vspace{1cm}
    {\large \Disciplina\par}
    \vspace{2.5cm}
    {\LARGE\bfseries \TituloRelatorio\par}
    \vspace{1.5cm}
    {\large Relatório acadêmico\par}
    \vfill
    \begin{minipage}{0.85\textwidth}
        \centering
        \textbf{Integrantes}\\[0.4cm]
        % Liste os integrantes em ordem alfabética, se essa for a orientação docente.
        Nome Sobrenome\\
        Nome Sobrenome\\
        Nome Sobrenome
    \end{minipage}
    \vfill
    {\large São José dos Campos\par}
    {\large \Ano\par}
\end{titlepage}

\tableofcontents
\newpage

\section{Introdução}
Apresente o problema estudado, sua motivação e a organização do relatório. Explique por que o método numérico escolhido é adequado ao problema.

\section{Objetivos}
\subsection{Objetivo geral}
Descreva o objetivo principal do trabalho.

\subsection{Objetivos específicos}
Liste os objetivos específicos, como implementar o método, testar sua precisão e analisar seu desempenho computacional.

\section{Fundamentação teórica}
Apresente os conceitos matemáticos necessários, as hipóteses e o método numérico utilizado. Descreva as etapas do algoritmo e cite as fontes consultadas.

\subsection{Descrição do método}
Explique o procedimento numérico. Inclua fórmulas, condições de aplicação, critérios de parada e possíveis limitações.

\subsection{Algoritmo}
Descreva as etapas do algoritmo em linguagem natural ou pseudocódigo. Se for útil, inclua o código-fonte em apêndice ou no repositório indicado na Seção \ref{sec:codigo}.

\section{Implementação computacional}
Informe a linguagem de programação, bibliotecas, ambiente de execução e decisões de implementação. Explique como o método foi traduzido para código e quais tolerâncias ou parâmetros foram utilizados.

\subsection{Trecho de código (opcional)}
\begin{lstlisting}[language=Python,caption={Exemplo de estrutura de código em Python}]
# Substitua este exemplo pelo código do trabalho
def metodo_numerico(dados):
    # implementar o método
    resultado = None
    return resultado
\end{lstlisting}

\section{Experimentos numéricos}
Descreva os problemas de teste, os dados utilizados, os valores iniciais, os parâmetros e as condições de execução. Os experimentos devem ser suficientemente documentados para permitir sua reprodução.

\subsection{Configuração dos testes}
Registre, quando aplicável, tolerância, número máximo de iterações, tamanho da malha, precisão aritmética, valores iniciais e critérios de parada.

\subsection{Resultados}
Apresente tabelas, gráficos e comparações. Informe as unidades e explique as métricas utilizadas.

\begin{table}[H]
\centering
\caption{Modelo de tabela de resultados}
\label{tab:resultados}
\begin{tabular}{lrrr}
\toprule
Problema & Resultado numérico & Erro estimado & Iterações\\
\midrule
Teste 1 & 0,0000 & 0,0000 & 0\\
Teste 2 & 0,0000 & 0,0000 & 0\\
\bottomrule
\end{tabular}
\end{table}

\subsection{Gráficos e figuras}
Envie os arquivos de imagem ao Overleaf e atualize o nome no comando abaixo.
\begin{figure}[H]
\centering
% Faça upload de figura.png para o projeto Overleaf.
\IfFileExists{figura.png}{%
    \includegraphics[width=0.75\textwidth]{figura.png}%
}{%
    \fbox{\parbox{0.7\textwidth}{\centering\vspace{1cm}Espaço reservado para figura (\texttt{figura.png})\vspace{1cm}}}%
}
\caption{Descrição do gráfico ou da figura.}
\label{fig:resultado}
\small Fonte: elaboração própria.
\end{figure}

\section{Análise e discussão}
Interprete os resultados. Discuta a precisão, a convergência, a estabilidade, o custo computacional e as limitações observadas, conforme o método estudado. Quando possível, compare com uma solução exata, analítica ou de referência.

\section{Conclusões}
Retome os objetivos e sintetize os principais resultados. Indique o que foi aprendido, as limitações do estudo e possíveis melhorias.

\section{Referências}
Liste todas as fontes citadas no texto. Cadastre as referências no arquivo \texttt{referencias.bib} e utilize comandos como \texttt{\textbackslash cite\{chave\}}.
\nocite{*}
\bibliography{referencias}

\appendix
\section{Código-fonte e instruções de execução}
\label{sec:codigo}
Inclua instruções para executar os programas. Para códigos extensos, prefira disponibilizar o material em um repositório e indicar o link abaixo.

\noindent\textbf{Repositório:} \url{https://github.com/organizacao/repositorio}

\section{Dados e problemas de teste}
Documente os arquivos de entrada, conjuntos de dados, exemplos, valores iniciais e parâmetros usados nos experimentos. Explique como reproduzir cada teste.

\section{Material complementar}
Inclua demonstrações, cálculos auxiliares, tabelas adicionais ou outras informações necessárias à compreensão do trabalho.

\end{document}
